
\label{app:gg_impl}
Our world model is initialized from BAGEL~\citep{deng2025emerging}, and largely uses the same training recipe. We use a subset of our robot data and egocentric human video data with high-quality segmented language labels, as we found that label quality (especially temporal segmentation quality) has a large impact on subgoal quality. We additionally mix in several open-source image editing datasets and open-source video datasets to better preserve the semantic knowledge of the model. Each training example consists of a subtask instruction, $\rawtext$, 3 camera inputs, $\bo_t$, and 3 target images, $\bo_{t_\text{end}}$, where $t_\text{end}$ is the last timestep of the segment spanning $t$. Following BAGEL, the camera inputs are processed using both a ViT (for semantic understanding) and VAE (for fine-grained image details). The ViT tokens are further processed by a 7B LLM backbone, and the VAE tokens are processed by a 7B generation backbone. ViT inputs are resized to a resolution of 448x336, and VAE inputs (including target images) are resized to a resolution of 512x384. The discrepancy is due to the differing patch size of the ViT and VAE (14 and 16, respectively). At test time, we set $\Delta = 4$ seconds as the time interval for regenerating subgoals to match SuSIE~\citep{black2023zero}.

\begin{figure*}[t]
\centering
\includegraphics[width=0.8\linewidth]{compressed_images_final/xemb_joint_vs_ee.pdf}\vspace{-5pt}
\caption{\textbf{Joint vs.\ end-effector control for prior models on cross-embodiment tasks.} We compare joint-space and end-effector (EE) control for baseline policies across a range of tasks, observing no substantial difference in performance between the two control modes.}
\label{fig:xemb_joint_vs_ee}
\end{figure*}

\begin{figure}[t]
    \centering
    \includegraphics[width=0.7\linewidth]{compressed_images_final/operator_experience_hours.pdf}\vspace{-5pt}
    \caption{\textbf{Operator experience in the human subject study.} Box plots show teleoperation experience (in hours) of the ten recruited operators across three categories: UR5e (target robot), the static bimanual robot (source robot), and all robots combined. The selected operators rank within the top 2\% of our operator fleet in terms of teleoperation experience.}
    \label{fig:operator_stats}
\end{figure}

\begin{figure}[t]
    \centering
    \includegraphics[width=0.7\linewidth]{compressed_images_final/human_vs_pi07.pdf}\vspace{-5pt}
    \caption{\textbf{Comparison of \MethodName{} (GC) and human.} We find that \MethodName{} (GC) achieves competitive performance compared to the human operators, in the shirt folding task with the UR5e bimanual platform.}
    \label{fig:human_vs_policy_quantitative}
\end{figure}


